% TOP_PROBLEM: {"id": 30006837, "title": "Global Regularity of the 2D Inviscid Boussinesq Equations", "category_id": 9, "category": {"id": 9, "name": "pde", "display_name": "Partial Differential Equations", "description": "PDEs and their applications in physics and geometry.", "slug": "pde", "order_index": 9, "created_at": "2026-07-31T15:26:25.671Z"}, "status": "open"}
% FIRST_POSTED: 2026-09-27
% !TeX program = lualatex
% ATTEMPT_STATUS: unresolved
\documentclass[11pt]{article}
\usepackage[margin=25mm]{geometry}
\usepackage{fontspec}
\setmainfont{Latin Modern Roman}
\usepackage{amsmath,amssymb,mathtools}
\usepackage{unicode-math}
\setmathfont{Latin Modern Math}
\usepackage{xurl,hyperref}
\hypersetup{colorlinks=true,urlcolor=blue}
\setcounter{secnumdepth}{0}
\setlength{\parindent}{0pt}
\setlength{\parskip}{5pt}
\setlength{\emergencystretch}{3em}
\begin{document}
\section*{\texorpdfstring{Global Regularity of the 2D Inviscid Boussinesq Equations}{Global Regularity of the 2D Inviscid Boussinesq Equations}}\label{30006837}
\subsection{Definitions and mathematical statement}
For $x\in{\mathbb{R}}^2$, the inviscid Boussinesq system in a standard normalized buoyancy convention is
\[
 \partial_tu+(u\cdot\nabla)u+\nabla p=\theta e_2,
 \qquad\partial_t\theta+u\cdot\nabla\theta=0,
 \qquad\nabla\cdot u=0.
\]
Here $u$ is velocity and $\theta$ the transported temperature or density. No viscosity or scalar diffusion is present. The question is whether every smooth finite-energy initial state in the cited Cauchy-problem class has global smooth evolution, or some such state loses regularity in finite time. The catalogue does not specify every Sobolev or decay norm hidden in its data description, so those hypotheses must follow the intended source. The domain is the whole plane; boundary or corner singularities do not by themselves settle this target.
\par\smallskip\textit{Formulation and concept references:} \href{https://arxiv.org/pdf/2306.08286}{[S1]}.

\subsection{Short English statement}
Can smooth finite-energy data for the undamped two-dimensional Boussinesq equations in the whole plane develop a finite-time singularity?
\subsection{Research attempt}
\paragraph{Scope and source check.}
This attempt was prepared by GPT-6 Astra Ultra on 27 September 2026.
Section 1.1 of [S1] discusses the undamped whole-plane problem separately
from the paper's partially dissipative theorems. The whole-plane
singularity result of Elgindi--Pasqualotto [E1] has $C^{1,\alpha}$
regularity with an exceptional radial smoothness point; it cannot be
silently substituted for the everywhere-smooth initial-data target.
Boundary singularities likewise have a different domain. No viscosity,
thermal diffusion, boundary, or external forcing is added here.

For precise a priori calculations take smooth rapidly decaying data,
or smooth data in $H^s(\mathbb R^2)$ for every sufficiently large $s$,
with divergence-free velocity. This is a concrete finite-energy
subclass of the intended smooth Cauchy problem. The work derives
energy and continuation estimates, gives an exact global example,
and disproves an overly strong vorticity-to-strain estimate. It does
not establish a global bound for every datum in the catalogue's class.

\paragraph{Transport and the energy bound.}
Along the smooth incompressible flow map, $\theta$ is constant and
Lebesgue measure is preserved. Therefore
\[
                  \|\theta(t)\|_{L^q}=\|\theta_0\|_{L^q}
                 \quad(1\le q\le\infty).               \tag{1}
\]
Taking the $L^2$ inner product of the velocity equation with $u$,
the transport and pressure terms integrate to zero, giving
\[
 \frac12\frac d{dt}\|u\|_2^2=\int\theta u_2
                      \le\|\theta_0\|_2\|u\|_2.
\]
Regularizing the norm by $\sqrt{\|u\|_2^2+\epsilon}$ and then taking
$\epsilon\downarrow0$ avoids division at a zero velocity norm. Thus
\[
                  \|u(t)\|_2\le\|u_0\|_2+t\|\theta_0\|_2.\tag{2}
\]
This controls kinetic energy on every finite interval of smooth
existence. It does not control the gradients that govern continuation.

\paragraph{The coupled derivative equations.}
Use the vorticity convention $\omega=\partial_1u_2-\partial_2u_1$.
Taking the curl of the momentum equation and differentiating the
transport equation gives
\[
 \begin{split}
 (\partial_t+u\cdot\nabla)\omega&=\partial_1\theta,\\
 (\partial_t+u\cdot\nabla)\nabla\theta
                           &=-(\nabla u)^T\nabla\theta.
 \end{split}                                             \tag{3}
\]
The sign of the first forcing term follows from the specified buoyancy
$\theta e_2$; changing that convention changes the curl sign.
Writing $X(t)=\|\nabla\theta(t)\|_\infty$,
$Y(t)=\|\omega(t)\|_\infty$, and $A(t)=\|\nabla u(t)\|_\infty$,
the flow formulas imply
\[
 X(t)\le X(0)\exp\left(\int_0^tA(s)\,ds\right),\qquad
 Y(t)\le Y(0)+\int_0^tX(s)\,ds.                         \tag{4}
\]
Conservation of $\|\theta\|_\infty$ in (1) does not conserve $X$.
The matrix in the second equation of (3) can stretch a gradient even
though the scalar values themselves are unchanged.

\paragraph{A conditional continuation estimate.}
For an integer $s\ge3$, the product and transport commutator estimates
give, with $F_s=\|u\|_{H^s}+\|\theta\|_{H^s}$,
\[
              \frac d{dt}F_s
                  \le C_s(1+A+X)F_s.                   \tag{5}
\]
To see the terms, commute a derivative $\partial^\alpha$ with
$u\cdot\nabla$. Its highest transport term integrates to zero by
$\nabla\cdot u=0$. The remaining products are bounded by
$\|\nabla u\|_\infty\|u\|_{H^s}$ in the velocity equation,
and by
$\|\nabla u\|_\infty\|\theta\|_{H^s}
 +\|\nabla\theta\|_\infty\|u\|_{H^s}$ in the scalar equation.
The forcing adds $\|\theta\|_{H^s}$; the pressure is removed by
the $L^2$-bounded Leray projection or by integration against
divergence-free differentiated velocity.

If $T<\infty$ and $\int_0^TA(t)\,dt<\infty$, then (4) bounds $X$
uniformly, and hence $\int_0^TX<\infty$. Gronwall in (5) bounds
$F_s$ on $[0,T)$. The local Sobolev continuation theorem for this
system, in the class described in [S1], then extends the solution.
Thus a finite-time smooth singularity in this class must have
\[
                              \int_0^TA(t)\,dt=\infty.   \tag{6}
\]
This is a conditional criterion, not a proof that its hypothesis
always holds. The derivative estimate and local continuation input
are separate parts of the argument.

\paragraph{An explicit obstruction to bounding strain by vorticity alone.}
A tempting closure is $A\le C Y$ with a universal constant. On the
whole plane it is false, even for smooth finite-energy velocities.
Choose a smooth radial cutoff $\chi_\epsilon$, bounded by one,
supported in $\epsilon<r<1$ and equal to one on
$2\epsilon\le r\le1/2$. Define
\[
                         \omega_\epsilon(r,\theta)
                                      =\chi_\epsilon(r)\sin2\theta.
\]
It is smooth, vanishes near the origin, has supremum norm at most one,
and has integral zero by angular symmetry. Let
\[
 u_\epsilon=K*\omega_\epsilon,\qquad
                   K(x)=\frac{(-x_2,x_1)}{2\pi|x|^2}.
\]
The zero integral makes the far-field velocity $O(|x|^{-2})$, by
subtracting $K(x)\int\omega_\epsilon$ and applying the mean value
bound to $K(x-y)-K(x)$. It is therefore in $L^2$; smoothness near
the compact vorticity support is standard convolution regularity,
or follows from the Poisson equation and smooth compact data.

There is no singularity in differentiating this convolution at zero,
because the vorticity vanishes near zero. Since
$\partial_1K_1(x)=\sin2\theta/(2\pi r^2)$,
\[
 \partial_1(u_\epsilon)_1(0)
   =\frac1{2\pi}\int\sin^2(2\theta)\,d\theta
                         \int_0^\infty\frac{\chi_\epsilon(r)}r\,dr
   =\frac12\int_0^\infty\frac{\chi_\epsilon(r)}r\,dr
   \ge\frac12\log\frac1{4\epsilon}.                     \tag{7}
\]
Thus $Y\le1$ while the velocity gradient can be arbitrarily large.
This is an exact family of spatial countertests to the proposed norm
inequality, not a time-dependent blow-up solution of Boussinesq.
Additional regularity norms are needed in a valid endpoint estimate.

Even the unjustified replacement $A=Y$ in the equality model
$Y'=X$, $X'=YX$ would not imply global boundedness. The explicit
functions $Y=2/(T-t)$, $X=2/(T-t)^2$ satisfy those ODEs and blow up.
An upper differential inequality permitting this behavior cannot by
itself exclude it. Conversely such ODE behavior does not force the PDE
to realize it, because the PDE has spatial constraints and signs.

\paragraph{A genuine smooth finite-energy global test.}
Set $r^2=x_1^2+x_2^2$ and
\[
 u(x)=e^{-r^2}(-x_2,x_1),\qquad
 \theta(x)=0,\qquad p(x)=-\tfrac14e^{-2r^2}.             \tag{8}
\]
The velocity is divergence free, smooth and rapidly decaying.
Its tangential speed is $v(r)=re^{-r^2}$, so its convective acceleration
is $-v(r)^2e_r/r=-re^{-2r^2}e_r$, canceled by
$\nabla p=re^{-2r^2}e_r$. Hence (8) is stationary for the full system.
Moreover
\[
                  \|u\|_2^2=2\pi\int_0^\infty r^3e^{-2r^2}\,dr
                              =\pi/4.
\]
This supplies an exact nonzero finite-energy velocity example.
It says nothing about all nonzero transported temperatures.
In contrast, a nonzero shear independent of one Cartesian coordinate
or a linear strain field on the entire plane generally has infinite
$L^2$ energy. Such ansatz solutions cannot be inserted into the
catalogue's finite-energy class without a new localization argument;
cutting them off changes the equations through pressure and transport.

\paragraph{Verification and final status.}
The local checks verify the curl convention, the radial pressure in
(8), the angular integral in (7), and the toy ODE identities. The
energy and continuation estimates are analytic proofs; numerical
smoothness for a finite time would not replace their missing global
bound. The lower-regularity source [E1] is cited with its actual scope.

\textbf{UNRESOLVED.} The complete partial results are (1)--(6), the
countertest (7), and the global example (8). The first missing step in
this proof route is an a priori integrable bound for $\|\nabla u\|_\infty$
for every allowed smooth datum, or an actual finite-energy smooth
construction violating it. Specific targets are a cancellation estimate
for the strain, control of the direction of $\nabla\theta$ along the
flow, and a localization mechanism for singularity models preserving
the precise whole-plane equation. No such result is proved here.

\subsection{Sources}
\begin{itemize}
\item[S1]Kyungkeun Kang, Jihoon Lee, and Dinh Duong Nguyen, Global well-posedness and stability of the 2D Boussinesq equations with partial dissipation (2024), Section 1.1.
\url{https://arxiv.org/pdf/2306.08286}.
\textit{Formulation source.}
\item[E1]Tarek M. Elgindi and Federico Pasqualotto, \emph{From Instability to Singularity Formation in Incompressible Fluids}.
\url{https://arxiv.org/abs/2310.19780}.
\textit{Whole-plane lower-regularity singularities; the exceptional smoothness hypothesis is not dropped.}
\end{itemize}
\end{document}
